% --- robocze: usunąć przed oddaniem pracy (razem z linią % --- robocze: usunąć przed oddaniem pracy (razem z linią % --- robocze: usunąć przed oddaniem pracy (razem z linią % --- robocze: usunąć przed oddaniem pracy (razem z linią \input{makra} w pliku głównym) ---
% Znacznik miejsca do weryfikacji, widoczny w PDF.
\newcommand{\dower}[1]{\textcolor{red}{[DO WERYFIKACJI: #1]}}
% Plan podrozdziału (z docs/02_spis_tresci.md), widoczny w PDF szarym drukiem.
\newenvironment{planroboczy}[1][]{%
  \par\begingroup\small\color{black!60}%
  \noindent\textbf{Plan roboczy\if\relax\detokenize{#1}\relax\else\ (#1)\fi}\par\nopagebreak
}{\par\endgroup\medskip}
\newcommand{\zrodlaplanowane}[1]{\par\noindent\textbf{Źródła planowane:} #1\par}
\newcommand{\materialyrobocze}[1]{\par\noindent\textbf{Materiały robocze:} #1\par}
 w pliku głównym) ---
% Znacznik miejsca do weryfikacji, widoczny w PDF.
\newcommand{\dower}[1]{\textcolor{red}{[DO WERYFIKACJI: #1]}}
% Plan podrozdziału (z docs/02_spis_tresci.md), widoczny w PDF szarym drukiem.
\newenvironment{planroboczy}[1][]{%
  \par\begingroup\small\color{black!60}%
  \noindent\textbf{Plan roboczy\if\relax\detokenize{#1}\relax\else\ (#1)\fi}\par\nopagebreak
}{\par\endgroup\medskip}
\newcommand{\zrodlaplanowane}[1]{\par\noindent\textbf{Źródła planowane:} #1\par}
\newcommand{\materialyrobocze}[1]{\par\noindent\textbf{Materiały robocze:} #1\par}
 w pliku głównym) ---
% Znacznik miejsca do weryfikacji, widoczny w PDF.
\newcommand{\dower}[1]{\textcolor{red}{[DO WERYFIKACJI: #1]}}
% Plan podrozdziału (z docs/02_spis_tresci.md), widoczny w PDF szarym drukiem.
\newenvironment{planroboczy}[1][]{%
  \par\begingroup\small\color{black!60}%
  \noindent\textbf{Plan roboczy\if\relax\detokenize{#1}\relax\else\ (#1)\fi}\par\nopagebreak
}{\par\endgroup\medskip}
\newcommand{\zrodlaplanowane}[1]{\par\noindent\textbf{Źródła planowane:} #1\par}
\newcommand{\materialyrobocze}[1]{\par\noindent\textbf{Materiały robocze:} #1\par}
 w pliku głównym) ---
% Znacznik miejsca do weryfikacji, widoczny w PDF.
\newcommand{\dower}[1]{\textcolor{red}{[DO WERYFIKACJI: #1]}}
% Plan podrozdziału (z docs/02_spis_tresci.md), widoczny w PDF szarym drukiem.
\newenvironment{planroboczy}[1][]{%
  \par\begingroup\small\color{black!60}%
  \noindent\textbf{Plan roboczy\if\relax\detokenize{#1}\relax\else\ (#1)\fi}\par\nopagebreak
}{\par\endgroup\medskip}
\newcommand{\zrodlaplanowane}[1]{\par\noindent\textbf{Źródła planowane:} #1\par}
\newcommand{\materialyrobocze}[1]{\par\noindent\textbf{Materiały robocze:} #1\par}
